% Paper 3, §8 — DeepSeek-R1 and the reasoning-trained model family
% Anchors: kb/notes/reasoning/reasoning-training.md §2-3
%          kb/notes/post-training/rlvr-and-grpo.md §2.3
%          kb/excerpts/deepseek-r1.md (#abstract, #sec-2, #sec-2-1-grpo,
%             #sec-2-1-hp, #sec-2-2-reward, #sec-2-3-aime, #sec-2-3-aha,
%             #sec-3-pipeline, #sec-4)
% Cross-paper: Paper-1 §long-context attention (precondition for 10K
%   traces); Paper-2 §SFT (cold-start and rejection-sampling SFT);
%   Paper-5 §contamination (AIME contamination contradiction).

\section{DeepSeek-R1 and the reasoning-trained model family}
\label{sec:r1-family}
% Alias label so the existing \cref{sec:r1-protocol} usages in §7 resolve
% to this section without requiring edits to that file.
\label{sec:r1-protocol}

If \cref{sec:inference-compute-scaling,sec:process-reward-models,sec:rlvr-grpo-dapo}
each treated one leg of the reasoning triangle in isolation, this
section is the empirical proof point that the three legs
\emph{compose}. \citet{deepseek-r1} train DeepSeek-R1 and DeepSeek-R1-Zero
on the DeepSeek-V3 base \citep[\S 2]{deepseek-r1}\footnote{KB excerpt:
\texttt{kb/excerpts/deepseek-r1.md\#sec-2}.} using a four-stage
pipeline whose RL phase moves the policy along the inference-compute
axis (response length grows from $\sim\!200$ to $\sim\!10^4$ tokens
\citep[\S 2.3]{deepseek-r1}\footnote{KB excerpt:
\texttt{kb/excerpts/deepseek-r1.md\#sec-2-3-aha}.}) using only the
verifiable terminal reward of \cref{eq:verifier-reward}, with no PRM
in the headline RL loop \citep[\S
2.2]{deepseek-r1}.\footnote{KB excerpt:
\texttt{kb/excerpts/deepseek-r1.md\#sec-2-2-reward}.} The headline
number --- AIME 2024 average pass@1 climbs from $15.6\%$ to $77.9\%$,
reaching $86.7\%$ with self-consistency over $64$ samples
\citep[\S 2.3]{deepseek-r1}\footnote{KB excerpt:
\texttt{kb/excerpts/deepseek-r1.md\#sec-2-3-aime}.} --- is the
cleanest single piece of evidence in the 2025--2026 reasoning literature
that the triangle is a co-design rather than three independent gains
that happen to ship together.

This section walks the four-stage R1 protocol, contrasts R1 with
R1-Zero (no cold start) to isolate what the SFT initialisation does
and does not contribute, transcribes the long-CoT compounding curve
that links response length to verifier accuracy, and traces the
distillation track and the open-weight reasoning ecosystem (Tülu~3,
Qwen3-thinking, QwQ, OpenR1) that R1 seeded. It closes on the
AIME-2024 contamination contradiction, the load-bearing open question
about whether the headline number is the gain it appears to be.

\subsection{The four-stage R1 protocol}
\label{sec:r1-fourstage}

\citet[\S 2.3]{deepseek-r1} establish a four-stage training protocol
for DeepSeek-R1 that subsequent open-replication work largely
follows.\footnote{KB note:
\texttt{kb/notes/reasoning/reasoning-training.md\#2}.} The pipeline is
the operational form of the reasoning triangle: each stage moves the
policy along a specific leg, and the legs compose in a fixed order.

\paragraph{Stage 1 --- cold-start SFT.} A small set
($\sim\!10^4$ examples) of curated long-CoT exemplars is used for SFT
on the DeepSeek-V3 base \citep[\S
3]{deepseek-r1}\footnote{KB excerpt:
\texttt{kb/excerpts/deepseek-r1.md\#sec-3-pipeline}.} The cold-start
data is described by \citet[\S 3]{deepseek-r1} as
``thousands of cold-start data that exhibits a conversational,
human-aligned thinking process.'' The SFT pass installs the
\texttt{<think>...</think><answer>...</answer>} format
\citep[\S 2]{deepseek-r1}\footnote{KB excerpt:
\texttt{kb/excerpts/deepseek-r1.md\#sec-2}.} that the format reward of
\cref{sec:rlvr-verifier-reward} subsequently maintains, and seeds the
language-consistency property that R1-Zero's purely-RL-trained traces
lack (\cref{sec:r1-zero-vs-r1}).

\paragraph{Stage 2 --- reasoning-oriented RL (GRPO with verifiable
rewards).} The cold-started policy is RL-trained with the GRPO
objective \cref{eq:grpo} of \cref{sec:rlvr-grpo-objective}, using the
binary verifier reward of \cref{eq:verifier-reward} on math, code,
and STEM problems with mechanically-checkable answers
\citep[\S 2.2]{deepseek-r1}.\footnote{KB excerpt:
\texttt{kb/excerpts/deepseek-r1.md\#sec-2-2-reward}.} The verifier is
a sum of an \emph{accuracy reward} (binary, from sympy-style
exact-match for math, unit-test execution for code, regex for
multiple-choice) and a small constant \emph{format reward} for the
\texttt{<think>...</think>} envelope
\citep[\S 2.2]{deepseek-r1}.\footnote{KB note:
\texttt{kb/notes/post-training/rlvr-and-grpo.md\#2-3}.} R1's protocol
explicitly excludes neural reward models: \citet[\S
2.2]{deepseek-r1} state ``we abstain from applying neural reward
models---whether outcome-based or process-based---to reasoning tasks''
on the grounds that learned RMs are reward-hackable under large-scale
RL.

The empirical content of Stage 2 is the headline of the entire R1
paper. Hyperparameters \citep[\S
2.1]{deepseek-r1}\footnote{KB excerpt:
\texttt{kb/excerpts/deepseek-r1.md\#sec-2-1-hp}.}: learning rate
$3\!\times\!10^{-6}$, KL coefficient $\beta = 0.001$, sampling
temperature $T = 1$ for rollout, group size $G = 16$, max trajectory
length $L_\text{max} = 32{,}768$ tokens before step $8.2\text{k}$ and
$65{,}536$ afterward, $10{,}400$ training steps total ($\approx 1.6$
epochs over the training distribution). Over this run, two coupled
quantities move together: average response length grows from a few
hundred tokens to $\sim\!10^4$, and AIME 2024 average pass@1 grows
from $15.6\%$ to $77.9\%$ \citep[\S
2.3]{deepseek-r1}.\footnote{KB excerpt:
\texttt{kb/excerpts/deepseek-r1.md\#sec-2-3-aime}.} \citet[\S
2.3]{deepseek-r1} additionally describe an
``aha moment'' --- a sudden spike in the policy's use of the
\texttt{wait} token during reflection passages.\footnote{KB excerpt:
\texttt{kb/excerpts/deepseek-r1.md\#sec-2-3-aha}.} The verifier is
agnostic to length and to reflection structure; the policy discovers
both as instrumentally useful for hitting the binary reward.

\paragraph{Stage 3 --- rejection-sampling SFT.} The Stage-2 RLVR
policy generates a large pool of solutions; the verifier filters for
correctness and a writing-quality filter further prunes the set,
yielding $\sim\!600$K reasoning traces that are concatenated with
$\sim\!200$K non-reasoning examples (general chat, safety formatting,
non-reasoning task prompts) into an $\sim\!800$K mixed
SFT set \citep[\S 3]{deepseek-r1}\footnote{KB excerpt:
\texttt{kb/excerpts/deepseek-r1.md\#sec-3-pipeline}.} The Stage-2
policy is fine-tuned on this set with standard cross-entropy SFT.
This stage is the bridge between ``specialist reasoner from Stage 2''
and ``general assistant that also reasons,'' and it is where the
reasoning capability becomes expressible as supervised data --- the
point that makes the distillation track of \cref{sec:r1-distillation}
possible at all.

\paragraph{Stage 4 --- general-domain RL.} A second RL phase combines
verifiable rewards (for the reasoning subset of prompts) with
preference-based rewards (for helpfulness, harmlessness, instruction
following on the broader prompt distribution) \citep[\S
3]{deepseek-r1}.\footnote{KB excerpt:
\texttt{kb/excerpts/deepseek-r1.md\#sec-3-pipeline}.} The released
DeepSeek-R1 weights are the output of this stage. Only Stage 2 of the
four is purely RLVR; Stage 4 is an RLVR/RLHF hybrid whose preference
component is the same machinery as the ``RL stage'' that bookends
post-training in Paper~2 §RL-stage. The two SFT stages (Stage 1 and
Stage 3) bookend the reasoning RL phase: cold-start SFT seeds it,
rejection-sampling SFT consolidates its outputs.

\subsection{R1-Zero versus R1: what does cold-start add?}
\label{sec:r1-zero-vs-r1}

DeepSeek-R1-Zero is the same Stage-2 RLVR run applied to the
DeepSeek-V3 base directly, without Stage 1's cold-start SFT
\citep[\S 2]{deepseek-r1}.\footnote{KB excerpt:
\texttt{kb/excerpts/deepseek-r1.md\#sec-2}.} The structural
comparison cleanly isolates what the cold-start contributes. R1-Zero
reaches \emph{the same reasoning capability headline} ---
$15.6\%\!\to\!77.9\%$ AIME 2024 pass@1, $86.7\%$ with
self-consistency \citep[\S 2.3]{deepseek-r1}\footnote{KB excerpt:
\texttt{kb/excerpts/deepseek-r1.md\#sec-2-3-aime}.} --- but its
traces are unreadable: ``DeepSeek-R1-Zero struggles with challenges
like poor readability, and language mixing, as DeepSeek-V3-Base is
trained on multiple languages, especially English and Chinese''
\citep[\S 3]{deepseek-r1}.\footnote{KB excerpt:
\texttt{kb/excerpts/deepseek-r1.md\#sec-3-pipeline}.} The traces
interleave Chinese and English at the token level; the format
envelope drifts; the prose is brittle. Capability and legibility
decouple.

The cold start, on this evidence, contributes three things
\citep[\S 3]{deepseek-r1}\footnote{KB note:
\texttt{kb/notes/reasoning/reasoning-training.md\#2}.}: \emph{format}
(the \texttt{<think>...</think><answer>...</answer>} envelope is
seeded so the format reward of \cref{sec:rlvr-verifier-reward} acts
on a non-zero baseline rather than carving the envelope from
scratch); \emph{language consistency} (a single language in the trace
because the cold-start exemplars are monolingual); and
\emph{regularisation} (a starting point inside the manifold of human
mathematical writing rather than at the multilingual base
distribution's mode). What it does \emph{not} contribute is reasoning
capability: that emerges in Stage 2 either way.

\begin{intuition}
R1-Zero is the pure-RL reasoning experiment; R1 is the same
experiment with a few thousand pages of human reasoning-style writing
mixed in at initialisation. Capability tracks the verifier signal;
legibility tracks what the policy was writing-like before that
signal arrived. Returning to math: GRPO's group-baseline
\cref{eq:grpo-advantage} only sees the binary verifier reward, so the
gradient is identical for a Chinese-English-mixed trace and a
clean-English trace at the same correctness label. The cold-start
biases the policy's initial $\pi_\theta$ toward one of those styles;
the RL phase preserves the bias because the gradient does not punish
it.
\end{intuition}

The R1-Zero / R1 split mirrors a long-standing pattern in RL: the
all-RL agent reaches comparable capability to the SFT-then-RL agent,
at the cost of some surface property the latter inherited from the
supervised initialisation.\footnote{KB note:
\texttt{kb/notes/reasoning/reasoning-training.md\#5} draws the
AlphaGo / AlphaGo-Zero analogy in this form.} The reading available
from R1's evidence: a small-but-curated supervised initialisation
substantially shortens the RL horizon to legibility without changing
the asymptote on capability.

\subsection{The long-CoT compounding curve}
\label{sec:r1-compounding}

The empirical signature of Stage 2 that links R1's evidence to the
inference-compute axis of \cref{sec:inference-compute-scaling} is the
\emph{long-CoT compounding effect}: response length and verifier
accuracy grow together over the RL run, in a relationship that is
qualitatively monotone-concave in
$\log\!L$.\footnote{KB note:
\texttt{kb/notes/reasoning/reasoning-training.md\#2-stage-2}; KB
excerpt: \texttt{kb/excerpts/deepseek-r1.md\#sec-2-3-aime} and
\texttt{\#sec-2-3-aha}.} Formally, let $L = |o|$ denote the
trajectory's token length and $p(L) = \E_{q, o \mid |o| \approx L}
[\mathbb{1}\{\mathrm{verify}(q, o) = \mathsf{true}\}]$ the verifier
pass rate at that length, both measured as the policy evolves
across training. \citet[\S 2.3]{deepseek-r1} report:
\begin{equation}
  p(L) \;\approx\; a\,\log L + b,
  \qquad L \in [\sim\!200,\;\sim\!10^4],
  \label{eq:r1-log-returns}
\end{equation}
qualitatively (no closed-form fit is reported in the paper; the
curve is shown as Figure~3 of \citet{deepseek-r1}). The same
qualitative shape is reported by \citet{s1-2025} on a small
reasoning-trained model under explicit budget forcing
(\cref{sec:budget-forcing}), and by the Snell-anchor empirics
\citep[\S 5--6]{snell2024}\footnote{KB excerpt:
\texttt{kb/excerpts/snell2024.md\#sec-5-2-search}.} on
test-time-strategy curves at fixed base models. As of 2026-Q2 no
parametric scaling law for $p(L)$ with the structural status of
Chinchilla's $L(N, D)$ has been attested in the literature; the
contradiction is catalogued at the paper level in
\cref{sec:rlvr-empirics} and revisited in §14.

\paragraph{Shape and tensor notes.} The training-time monotone-concave
curve \cref{eq:r1-log-returns} is a \emph{joint} property of the
policy and the verifier: the policy emits longer traces because longer
traces hit the binary reward more often on the prompts the dataset
samples; the verifier supplies the signal that makes longer traces
cost-justified. Trace-length growth is therefore not a separate
training objective with its own loss term --- the GRPO objective
\cref{eq:grpo} contains no length term --- but a fixed-point of the
RL dynamics under sparse verifiable reward
\citep[\S 2.2.2]{deepseek-r1}.\footnote{KB note:
\texttt{kb/notes/reasoning/reasoning-training.md\#5}.} The shape of
the curve at inference time, holding the policy fixed and varying
$L$ via budget forcing, is the subject of \cref{sec:budget-forcing};
the shape during training, with the policy itself moving along $L$,
is the subject of this section.

\begin{analogy}
The trace-length growth of Stage 2 is a policy-improvement fixed
point under sparse reward: longer trace $\Rightarrow$ more chances to
self-correct $\Rightarrow$ higher verifier pass rate $\Rightarrow$
GRPO gradient prefers the longer-trace mode of the rollout
distribution $\Rightarrow$ next iteration's expected length increases.
Returning to canonical form: in
\cref{eq:grpo-advantage} the trajectory-level $\tilde{A}_i$ is
positive on the $i$-th rollout exactly when $R_i$ is above the
group's mean; on prompts where the long-trace rollouts hit
$R_i = 1$ more often than the short-trace ones, those long-trace
rollouts carry positive standardised advantage and pull the policy's
expected length upward through repeated iterations of \cref{eq:grpo}.
The dynamics are the verifier-conditioned analogue of standard
policy-gradient ascent on a reward whose support is a subset of
``what the policy already does sometimes'' --- the same mechanism
flagged in \cref{sec:rlvr-grpo-objective}'s intuition box about
RLVR's reach being bounded by the base policy's pass@$K$ at any $K$
\citep{rlvr-limits-2025}.
\end{analogy}

The long-context architectural prerequisite is non-negotiable. Trace
lengths in $[10^4, 6.5\!\times\!10^4]$ tokens
\citep[\S 2.1]{deepseek-r1} require an attention substrate that
costs sub-quadratically per token --- some combination of MLA, GQA,
or sliding-window attention --- to make the per-step decoder
inference cost tractable at training-time rollout volume. Paper~1
§long-context-attention is the precondition for this section's
$10^4$-token traces; without the architectural innovations of
DeepSeek-V2/V3 there is no policy that can be RL-trained on
trajectories of this length within feasible compute.

\subsection{The distillation track and the open-weight ecosystem}
\label{sec:r1-distillation}

\citet[\S 4]{deepseek-r1}\footnote{KB excerpt:
\texttt{kb/excerpts/deepseek-r1.md\#sec-4}.} include a separate track
that distills R1 into smaller dense bases via \emph{pure SFT on the
800K Stage-3 set} (the same $600$K reasoning + $200$K non-reasoning
data of \cref{sec:r1-fourstage}). The distillation targets span the
$1.5$B--$70$B range using Qwen2.5 and LLaMA-3 bases. The notable
empirical result: the distilled $32$B model matches or exceeds
OpenAI's o1-mini on multiple reasoning benchmarks, and the distilled
$7$B is competitive with strong baselines at $7$B scale, all without
any RL on the small models themselves \citep[\S 4]{deepseek-r1}. The
implication \citet{deepseek-r1} highlight is structural: \emph{the
reasoning capability that Stage 2 elicits is expressible as a (large)
SFT dataset}. The small-model trainer does not need to run RLVR; the
trainer at the teacher scale does, but only once.

\begin{intuition}
At small scales the bottleneck on long-CoT capability is data ---
trace exemplars at the right depth and difficulty --- not algorithm.
Once a teacher emits a sufficient corpus of correct long-CoT traces
verified against ground truth, vanilla SFT on those traces transfers
the capability efficiently. RLVR's role in the lineage is to
\emph{produce} that corpus when no prior trace-generating
distribution exists; the RL is needed once at teacher scale, not at
every student scale. Returning to math: the SFT cross-entropy on a
verifier-filtered Stage-3-style trace set is an upper bound on the
KL divergence the student would need to traverse to reach the
teacher's verifier-conditional distribution; once the bound is small
enough, the student inherits the teacher's pass-rate without
on-policy verifier signal.
\end{intuition}

The distillation track seeded a measurable open-weight reasoning
ecosystem in 2025--2026, of which the most visible entries are:
\textbf{Tülu~3} \citep{tulu3} --- a reproducible RLVR-on-Llama
recipe with verifier rewards on math/code prompts, released alongside
training data; \textbf{OpenR1}
\deepencite{openr1 — citation-pending} --- a Hugging Face community
project that reproduced R1-Zero-style runs on smaller bases with
public code; \textbf{Qwen3-thinking}
\deepencite{qwen3-thinking — citation-pending; cf.\ \citet{qwen3} for
the broader Qwen3 lineage} --- a single model that toggles between
thinking and non-thinking modes via system prompt, with mode-tagged
post-training; \textbf{QwQ}
\deepencite{qwq — citation-pending} --- Qwen's standalone reasoning
preview that predates Qwen3-thinking. The structural commitment
shared across this ecosystem is the R1 protocol's order: cold-start
SFT, then RLVR on verifiable rewards, then either rejection-sampling
SFT (open-weight reproductions) or its functional equivalent
(commercial closed-weight reasoning models). \cref{sec:frontier-2026}
returns to which entries in the ecosystem ship which legs of the
triangle and at what budgets.

\subsection{The AIME-2024 contamination contradiction}
\label{sec:r1-contamination}

The R1 headline number is the cleanest single piece of empirical
evidence for the reasoning triangle's composability, which makes its
contamination status the load-bearing open question of this section.

\begin{contradiction}
\textbf{Was a fraction of AIME-2024-style problems plausibly in
DeepSeek-V3's pretraining corpus?} The $15.6\%\!\to\!77.9\%$ AIME 2024
pass@1 jump \citep[\S 2.3]{deepseek-r1} attributes the gain to
RLVR-elicited reasoning. The competing reading is that AIME-style
problems and their solutions are widely indexed on the open web ---
in problem-set archives, contest-prep textbooks, math forums, and
AoPS-style community posts --- and that \emph{some fraction} of the
pretraining-corpus distribution overlaps with the test
distribution. Under that reading, part of the headline gain is
recall of memorised solutions surfaced through long-CoT
self-prompting, not novel reasoning composed under RL. Reported
contamination-detection methodologies (n-gram overlap on solutions,
near-duplicate matching, decontamination-by-substring) all have
known false-negative modes against post-hoc-paraphrased problems and
arithmetic isomorphisms (date shifts, variable renames). \citet{deepseek-r1}
do not report a contamination audit on the AIME 2024 train/test
boundary, and as of 2026-Q2 no third-party study has reproduced an
audit on R1's pretraining mix. The cross-paper treatment of
contamination methodology is Paper~5 §contamination, which catalogues
the methodologies that disagree on what counts as contamination, and
Paper~3 §14 (contradictions) returns to this specific
AIME-2024-on-R1 question with the post-2026 evidence base.
\end{contradiction}

The contamination contradiction interacts non-trivially with the
\citet{rlvr-limits-2025} ``no new capabilities'' result of
\cref{sec:rlvr-grpo-objective} and the contradictions of §14. If the
base model's pass@$K$ at moderate $K$ on AIME-style problems is
already high because of pretraining contamination, then the
\citeauthor{rlvr-limits-2025} reading --- RLVR sharpens pass@1 but
does not expand the solvable set --- becomes near-tautological on
this benchmark. The R1 evidence does not distinguish ``RL elicited
reasoning capability that the base model could not realise on its
own at any sample budget'' from ``RL surfaced memorised solution
patterns the base model held latently from pretraining.'' Both
readings are compatible with the headline curve; only a
contamination-controlled rerun on a held-out reasoning benchmark
would discriminate.

\subsection{Forward link}
\label{sec:r1-forward}

This section established the empirical proof point that the
reasoning triangle composes: the four-stage R1 protocol moves the
DeepSeek-V3 base from non-reasoning capability to AIME-2024 pass@1
of $77.9\%$ (and $86.7\%$ with self-consistency) using GRPO
\cref{eq:grpo} on the verifier reward \cref{eq:verifier-reward}, with
no PRM in the headline RL loop and no neural reward model anywhere
in the pipeline \citep[\S 2.2]{deepseek-r1}. The R1-Zero comparison
isolates cold-start SFT's role to legibility and format rather than
capability; the long-CoT compounding curve \cref{eq:r1-log-returns}
links the response-length growth observed during training to the
inference-compute axis of \cref{sec:inference-compute-scaling}; the
distillation track \citep[\S 4]{deepseek-r1} demonstrates that the
elicited reasoning capability is expressible as a supervised dataset
once.

Three downstream sections pick up directly from R1's evidence.
\Cref{sec:budget-forcing} generalises the s1 \citep{s1-2025} minimal
recipe --- a small model, $1$K reasoning SFT examples, and explicit
budget-forcing of the \texttt{<think>} channel --- as the
inference-time complement to R1's training-time long-CoT discovery,
and tests whether the same monotone-concave-in-$\log L$ shape that
\cref{eq:r1-log-returns} observed during training holds at inference
on a fixed reasoning-trained policy.
\Cref{sec:cot-faithfulness} opens the question whether R1's traces
causally drive the answers --- whether the verifier's binary signal
selects for traces $z$ such that
$p_\theta(y \mid x, z) \neq p_\theta(y \mid x, z')$ on
intervention-corrupted $z'$ --- which is the load-bearing question
for any downstream use of CoT as a probe surface (Paper~4
§interpretability) or a monitorable channel (Paper~5 §alignment).
\Cref{sec:prm-survival} re-asks the
contradiction of \cref{sec:prm-post-r1}: given that R1 achieved its
headline with terminal verifier rewards alone, do PRMs survive the
reasoning-trained-model era as a training signal, or only as
inference-time rerankers? The R1 evidence is consistent with either
reading; the distinguishing empirics live in the post-R1 PRM
literature catalogued in \cref{sec:prm-survival}.
